\documentclass[14pt,a4paper,oneside]{extbook}
\usepackage[utf8]{vietnam}
\usepackage{amsmath,amssymb,mathrsfs,fancyhdr,enumerate,multirow,makecell}
\usepackage{tkz-euclide,tikz-3dplot,tikz,tkz-tab}
\usetikzlibrary{arrows.meta,calc,intersections,angles,quotes,backgrounds,shapes.geometric,patterns,through,shadows,positioning}
\usepackage{yhmath}
\usepackage[top=0.7cm, bottom=0.7cm, left=1.0cm, right=1.0cm, headheight=16pt, headsep=3mm, footskip=5mm]{geometry}
\usepackage[hidelinks,unicode]{hyperref}
\usepackage{qrcode}
\usepackage[many]{tcolorbox}
\usepackage{fontawesome}

\pagestyle{fancy}
\fancyhf{}
\fancyhead[L]{\footnotesize\bfseries\color{teal!85!black}\faBook~SỔ TAY KIẾN THỨC TOÁN 5}
\fancyhead[R]{\footnotesize\bfseries\color{blue!85!black}Lớp Toán Thầy Lĩnh -- loptoanthaylinh.vn}
\fancyfoot[C]{\footnotesize\bfseries\color{gray!80!black}-- Sổ tay kiến thức: Bài 10. Khái niệm số thập phân --}
\renewcommand{\headrulewidth}{0.4pt}

\begin{document}

% --- BANNER TIÊU ĐỀ BÀI HỌC ---
\noindent
\begin{tcolorbox}[
  enhanced,
  colback=teal!8!white,
  colframe=teal!80!black,
  arc=3mm,
  boxrule=1.1pt,
  width=\linewidth,
  halign=center,
  top=2.2mm, bottom=2.2mm,
  left=3mm, right=3mm,
  drop fuzzy shadow
]
  {\large\bfseries\color{teal!95!black}\faBookmark\ BÀI 10. KHÁI NIỆM SỐ THẬP PHÂN}\\[2pt]
  {\scriptsize\color{black!75}\itshape Cấu tạo phần nguyên, phần thập phân, các hàng thập phân và cách đọc viết trên 1 trang A4}
\end{tcolorbox}

\vspace{3.2mm}
% --- KHỐI 1: CẤU TẠO SỐ THẬP PHÂN ---
\noindent
\begin{minipage}[t]{0.58\textwidth}
  \begin{tcolorbox}[
    enhanced,
    title={\footnotesize\bfseries\color{white}\faKey\ 1. Cấu tạo của số thập phân},
    colback=blue!2!white,
    colframe=blue!75!black,
    arc=2.5mm, boxrule=0.8pt, top=1.8mm, bottom=1.8mm, left=2.2mm, right=2.2mm
  ]
    \footnotesize
    Mỗi số thập phân gồm hai phần cách nhau bởi \textbf{dấu phẩy}:\\
    $\bullet$ \textbf{Phần nguyên:} Đứng bên \textbf{trái} dấu phẩy (như số tự nhiên: đơn vị, chục, trăm, nghìn\dots).\\
    $\bullet$ \textbf{Phần thập phân:} Đứng bên \textbf{phải} dấu phẩy.\\
    $\triangleright$ Chữ số đầu tiên sau dấu phẩy là hàng \textbf{phần mười} ($\frac{1}{10}$).\\
    $\triangleright$ Chữ số thứ hai là hàng \textbf{phần trăm} ($\frac{1}{100}$).\\
    $\triangleright$ Chữ số thứ ba là hàng \textbf{phần nghìn} ($\frac{1}{1000}$).
  \end{tcolorbox}
\end{minipage}%
\hfill
\begin{minipage}[t]{0.40\textwidth}
  \begin{tcolorbox}[
    enhanced,
    title={\footnotesize\bfseries\color{white}\faEye\ Ví dụ phân tích số $35,428$},
    colback=teal!2!white,
    colframe=teal!75!black,
    arc=2.5mm, boxrule=0.8pt, top=1.8mm, bottom=1.8mm, left=1.5mm, right=1.5mm,
    halign=center
  ]
    \scriptsize
    \begin{tabular}{|c|c|c|c|c|}
      \hline
      \multicolumn{2}{|c|}{\textbf{Phần nguyên}} & \multicolumn{3}{c|}{\textbf{Phần thập phân}} \\ \hline
      Chục & Đ.vị & $\frac{1}{10}$ & $\frac{1}{100}$ & $\frac{1}{1000}$ \\ \hline
      $3$ & $5$ & $4$ & $2$ & $8$ \\ \hline
    \end{tabular}\\[3pt]
    {\tiny\color{teal!95!black}\bfseries Giá trị của $4$ là $\frac{4}{10}$; của $2$ là $\frac{2}{100}$; của $8$ là $\frac{8}{1000}$}
  \end{tcolorbox}
\end{minipage}

\vspace{3mm}

% --- KHỐI 2: ĐỌC VIẾT & SỐ THẬP PHÂN BẰNG NHAU ---
\noindent
\begin{minipage}[t]{0.50\textwidth}
  \begin{tcolorbox}[
    enhanced,
    title={\footnotesize\bfseries\color{white}\faPencil\ 2. Đọc và viết số thập phân},
    colback=purple!2!white,
    colframe=purple!70!black,
    arc=2.5mm, boxrule=0.8pt, top=1.8mm, bottom=1.8mm, left=2.2mm, right=2.2mm
  ]
    \footnotesize
    $\bullet$ Đọc phần nguyên trước, đọc dấu ``\textbf{phẩy}'', rồi đọc phần thập phân.\\
    \textit{VD:} $8,45$: Tám phẩy bốn mươi lăm.\\
    $\bullet$ Viết lần lượt từ hàng cao xuống hàng thấp từ trái qua phải.
  \end{tcolorbox}
\end{minipage}%
\hfill
\begin{minipage}[t]{0.48\textwidth}
  \begin{tcolorbox}[
    enhanced,
    title={\footnotesize\bfseries\color{white}\faEqual\ 3. Số thập phân bằng nhau},
    colback=orange!2!white,
    colframe=orange!80!black,
    arc=2.5mm, boxrule=0.8pt, top=1.8mm, bottom=1.8mm, left=2mm, right=2mm
  ]
    \footnotesize
    $\bullet$ Nếu viết thêm chữ số $0$ vào \textbf{bên phải tận cùng} của phần thập phân thì được số thập phân \textbf{bằng nó}.\\
    \textit{VD:} $0,9 = 0,90 = 0,900$.\\
    $\bullet$ Nếu bớt đi các chữ số $0$ ở tận cùng bên phải phần thập phân thì giá trị \textbf{không đổi}.
  \end{tcolorbox}
\end{minipage}

\vspace{3mm}

% --- KHỐI 3: SƠ ĐỒ TƯ DUY ---
\noindent
\begin{tcolorbox}[
  enhanced,
  title={\footnotesize\bfseries\color{white}\faSitemap\ Sơ đồ tư duy: Các hàng của số thập phân},
  colback=yellow!2!white,
  colframe=teal!85!black,
  arc=3mm, boxrule=0.9pt, top=1.8mm, bottom=1.8mm, left=1.5mm, right=1.5mm,
  halign=center
]
  \begin{tikzpicture}[
    scale=0.92,
    every node/.style={transform shape},
    c0/.style={rectangle, rounded corners=5pt, fill=teal!85!black, text=white, font=\bfseries\footnotesize, inner sep=5pt, align=center},
    c1/.style={rectangle, rounded corners=4pt, fill=blue!15, draw=blue!70!black, text=blue!90!black, font=\bfseries\scriptsize, inner sep=4.5pt, align=center},
    c2/.style={rectangle, rounded corners=4pt, fill=orange!15, draw=orange!70!black, text=orange!90!black, font=\bfseries\scriptsize, inner sep=4.5pt, align=center}
  ]
    \node[c0] (core) at (0, 0.75) {\faTag\ DẤU PHẨY ( , )};
    \node[c1] (left) at (-4.7, -0.45) {\textbf{BÊN TRÁI: Phần nguyên}\\[2pt](Trăm $\leftarrow$ Chục $\leftarrow$ Đơn vị)};
    \node[c2] (right) at (4.7, -0.45) {\textbf{BÊN PHẢI: Phần thập phân}\\[2pt](Phần mười $\rightarrow$ Phần trăm $\rightarrow$ Phần nghìn)};

    \draw[teal!85!black, thick, -{Stealth[length=2mm]}] (core) -| (left);
    \draw[teal!85!black, thick, -{Stealth[length=2mm]}] (core) -| (right);
  \end{tikzpicture}
\end{tcolorbox}

\vspace{3mm}

% --- KHỐI 4: CHUYỂN ĐỔI PSTP & BẪY TRẮC NGHIỆM ---
\noindent
\begin{minipage}[t]{0.54\textwidth}
  \begin{tcolorbox}[
    enhanced,
    title={\footnotesize\bfseries\color{white}\faBolt\ Đổi nhanh Phân số thập phân $\leftrightarrow$ STP},
    colback=green!3!white,
    colframe=green!60!black,
    arc=2.5mm, boxrule=0.8pt, top=1.8mm, bottom=1.8mm, left=2.2mm, right=2.2mm
  ]
    \footnotesize
    Số chữ số sau dấu phẩy bằng đúng \textbf{số chữ số 0} của mẫu:\\
    $\bullet$ $\dfrac{7}{10} = 0,7$ ($1$ chữ số $0 \rightarrow 1$ chữ số sau phẩy).\\
    $\bullet$ $\dfrac{3}{100} = 0,03$ ($2$ chữ số $0 \rightarrow 2$ chữ số sau phẩy).\\
    $\bullet$ $\dfrac{125}{100} = 1,25$ \qquad $\bullet$ $\dfrac{9}{1000} = 0,009$.
  \end{tcolorbox}
\end{minipage}%
\hfill
\begin{minipage}[t]{0.44\textwidth}
  \begin{tcolorbox}[
    enhanced,
    title={\footnotesize\bfseries\color{white}\faExclamationTriangle\ Mẹo tránh bẫy trắc nghiệm},
    colback=red!2!white,
    colframe=red!70!black,
    arc=2.5mm, boxrule=0.8pt, top=1.8mm, bottom=1.8mm, left=2.2mm, right=2.2mm
  ]
    \footnotesize
    $\bullet$ Học sinh rất hay viết $\dfrac{3}{100} = 0,3$ (\textbf{SAI}). Phải là $0,03$.\\
    $\bullet$ Số $0$ ở giữa phần thập phân \textbf{không được bỏ}: $5,06 \neq 5,6$ (chỉ được bỏ số $0$ ở tận cùng bên phải!).
  \end{tcolorbox}
\end{minipage}

\end{document}
